\section*{अध्याय \ref{ch:g8:winds-global-circulation} --- पवनें और वैश्विक परिसंचरण}

\begin{solution}{exo:g8:winds-global-circulation:1}
(क) उत्तर-पूर्वी व्यापारिक पवन, उत्तर-पूर्व से। (ख) \omterm{def:g8:winds-global-circulation:global-winds}{पछुआ पवनें}, उत्तर-पश्चिम से।
(ग) \omterm{def:g8:winds-global-circulation:global-winds}{ध्रुवीय पुरवा पवनें}, पूर्व से (उत्तर-पूर्व से)। (घ) दक्षिण-पूर्वी व्यापारिक पवन,
दक्षिण-पूर्व से।
\end{solution}

\begin{solution}{exo:g8:winds-global-circulation:2}
(क) गरम हवा ऊपर उठती है और ठंडी हवा नीचे उतरती है। (ख) निम्न। (ग) उच्च दाब से निम्न
दाब की ओर। (घ) दाईं ओर।
\end{solution}

\begin{solution}{exo:g8:winds-global-circulation:3}
(क) ग़लत: \omterm{def:g8:winds-global-circulation:breezes}{समुद्र समीर} \omterm{def:g2:water-on-earth:ocean}{समुद्र} से ज़मीन की ओर बहता है (दिन में)। (ख) सही। (ग) ग़लत:
वे पूर्व से बहती हैं (उत्तर-पूर्व से या दक्षिण-पूर्व से)। (घ) सही।
\end{solution}

\begin{solution}{exo:g8:winds-global-circulation:4}
चित्र एक चक्कर दिखाता है: ज़मीन के साथ-साथ \omterm{def:g2:water-on-earth:ocean}{समुद्र} से ज़मीन की ओर एक तीर; ज़मीन के
ऊपर, ऊपर जाता एक तीर; ऊँचाई पर, ज़मीन से वापस \omterm{def:g2:water-on-earth:ocean}{समुद्र} के ऊपर एक तीर; \omterm{def:g2:water-on-earth:ocean}{समुद्र} के ऊपर,
नीचे आता एक तीर। उतरती हवा के नीचे (\omterm{def:g2:water-on-earth:ocean}{समुद्र} के ऊपर) ज़मीन पर ``उच्च'' लिखा है, उठती
हवा के नीचे (ज़मीन के ऊपर) ``निम्न''।
\end{solution}

\begin{solution}{exo:g8:winds-global-circulation:5}
(क) 1.41 गुना ज़मीन। (ख) 50~\%। (ग) 26~\%। (घ) \omterm{def:g3:continents-and-oceans:poles}{ध्रुवों} के पास सूर्य का प्रकाश
तिरछा आता है और एक बड़े क्षेत्र पर फैल जाता है, इसलिए हर वर्ग मीटर को थोड़ी ऊर्जा
मिलती है।
\end{solution}

\begin{solution}{exo:g8:winds-global-circulation:6}
(क) \omterm{def:g8:winds-global-circulation:global-winds}{ध्रुवीय पुरवा पवनें} (\ang{70} से \ang{60} उत्तर तक), \omterm{def:g8:winds-global-circulation:global-winds}{पछुआ पवनें} (\ang{60} से
\ang{30} उत्तर तक), उत्तर-पूर्वी \omterm{def:g8:winds-global-circulation:global-winds}{व्यापारिक पवनें} (\ang{30} उत्तर से \omterm{def:g3:continents-and-oceans:equator}{भूमध्य रेखा} तक),
दक्षिण-पूर्वी \omterm{def:g8:winds-global-circulation:global-winds}{व्यापारिक पवनें} (\omterm{def:g3:continents-and-oceans:equator}{भूमध्य रेखा} से \ang{30} दक्षिण तक), \omterm{def:g8:winds-global-circulation:global-winds}{पछुआ पवनें}
(\ang{30} से \ang{60} दक्षिण तक), \omterm{def:g8:winds-global-circulation:global-winds}{ध्रुवीय पुरवा पवनें} (\ang{60} से \ang{70}
दक्षिण तक)। (ख) \omterm{def:g3:continents-and-oceans:equator}{भूमध्य रेखा} के पास, डोलड्रम में। (ग) \ang{30} उत्तर और
\ang{30} दक्षिण के पास, नीचे उतरती हवा के नीचे।
\end{solution}

\begin{solution}{exo:g8:winds-global-circulation:7}
(क) ज़मीन ($31 > 22$): पवन \omterm{def:g2:water-on-earth:ocean}{समुद्र} से ज़मीन की ओर बहती है, एक \omterm{def:g8:winds-global-circulation:breezes}{समुद्र समीर}।
(ख) \omterm{def:g2:water-on-earth:ocean}{समुद्र} ($21 > 17$): पवन ज़मीन से \omterm{def:g2:water-on-earth:ocean}{समुद्र} की ओर बहती है, एक \omterm{def:g8:winds-global-circulation:breezes}{स्थल समीर}।
(ग) ज़मीन: $31 - 17 = 14$ डिग्री; \omterm{def:g2:water-on-earth:ocean}{समुद्र}: $22 -
21 = 1$ डिग्री। (घ) सूर्य की किरणें
ज़मीन की केवल एक पतली परत को गरम करती हैं, जो जल्दी गरम और जल्दी ठंडी होती है; \omterm{def:g2:water-on-earth:ocean}{समुद्र}
अपनी ऊष्मा पानी की एक मोटी परत में फैला देता है, इसलिए उसका तापमान दिन और रात के बीच
मुश्किल से बदलता है।
\end{solution}

\begin{solution}{exo:g8:winds-global-circulation:8}
(क) उत्तर से। (ख) दाईं ओर। (ग) उत्तर-पूर्व से (यह उत्तर-पूर्वी व्यापारिक पवन है)।
(घ) घूमने के बिना वह दक्षिण से बहती; घूमना उसे बाईं ओर मोड़ता है; वह दक्षिण-पूर्व से
बहती है।
\end{solution}

\begin{solution}{exo:g8:winds-global-circulation:9}
\omterm{def:g3:continents-and-oceans:equator}{भूमध्य रेखा} के पास ज़मीन और \omterm{def:g2:water-on-earth:ocean}{समुद्र} सूर्य से बहुत \emph{गरम होते हैं}; गरम, नम हवा
\emph{ऊपर उठती है}; उठते हुए वह \emph{ठंडी होती है}, और उसकी जलवाष्प
\emph{\omterm{def:g4:clouds-rain-snow:cloud}{बादलों}} में बदल जाती है, जो बौछारें और तड़ित-झंझा लाते हैं। \ang{25} उत्तर के
पास, \omterm{def:g3:continents-and-oceans:equator}{भूमध्य रेखा} पर उठी हवा \emph{नीचे उतरती है}; उसके नीचे ज़मीन पर दाब
\emph{उच्च} होता है; उतरती हवा गरम होती है, उसके \omterm{def:g4:clouds-rain-snow:cloud}{बादल} सूख जाते हैं और क्षेत्र
\emph{सूखा} रहता है: एक \omterm{def:g5:weather-and-climate:desert}{मरुस्थल}।
\end{solution}

\begin{solution}{exo:g8:winds-global-circulation:10}
(क) \omterm{def:g8:winds-global-circulation:global-winds}{पछुआ पवनें}। नहीं: वे पश्चिम से बहती हैं, पश्चिम की ओर चलते जहाज़ के ख़िलाफ़।
(ख) \omterm{def:g8:winds-global-circulation:global-winds}{व्यापारिक पवनों} की पेटी तक पहुँचने के लिए, लगभग \ang{30} उत्तर के दक्षिण में।
(ग) उत्तर-पूर्वी व्यापारिक पवन, \omterm{def:g3:continents-and-oceans:equator}{भूमध्य रेखा} और लगभग \ang{30} उत्तर के बीच।
(घ) \omterm{def:g8:winds-global-circulation:global-winds}{पछुआ पवनें}, लगभग \ang{30} उत्तर के उत्तर में (\ang{40} उत्तर के पास, जो घर के
बंदरगाह का \omterm{def:g6:latitude-longitude:latitude}{अक्षांश} है)।
\end{solution}

\begin{solution}{exo:g8:winds-global-circulation:11}
(क) $1 \div 0.5 = 2$: तालिका 2 देती है। (ख) $1 \div 0.866 \approx
1.1547$, यानी
\textbf{1.15}। (ग) \ang{30} पर: $1000 \div 1.15 \approx
\qty{870}{W}$ हर वर्ग
मीटर पर (ठीक-ठीक $1000 \times 0.866 =
\qty{866}{W}$)। \ang{60} पर:
$1000 \div 2 = \qty{500}{W}$।
\end{solution}

\begin{solution}{exo:g8:winds-global-circulation:12}
(क) तीन चक्कर: हवा \omterm{def:g3:continents-and-oceans:equator}{भूमध्य रेखा} पर उठती है और \ang{30} के पास उतरती है; \ang{60} के
पास उठती है और \ang{30} के पास उतरती है; \ang{60} के पास उठती है और \omterm{def:g3:continents-and-oceans:poles}{ध्रुव} पर उतरती
है। (ख) \omterm{def:g3:continents-and-oceans:equator}{भूमध्य रेखा} पर और \ang{60} के पास निम्न; \ang{30} के पास और \omterm{def:g3:continents-and-oceans:poles}{ध्रुव} पर उच्च।
(ग) \ang{30} से \omterm{def:g3:continents-and-oceans:equator}{भूमध्य रेखा} तक ज़मीनी हवा दक्षिण की ओर जाती है और दाईं ओर मुड़ती है:
उत्तर-पूर्वी व्यापारिक पवन। \ang{30} से \ang{60} तक वह उत्तर की ओर जाती है और दाईं ओर
मुड़ती है: वह दक्षिण-पश्चिम से आती है। \omterm{def:g3:continents-and-oceans:poles}{ध्रुव} से \ang{60} तक वह दक्षिण की ओर जाती है,
दाईं ओर मुड़ी हुई: ध्रुवीय पुरवा पवन। (घ) बीच वाले कोष्ठ में ज़मीनी हवा \omterm{def:g3:continents-and-oceans:poles}{ध्रुव} की ओर,
यानी उत्तर की ओर, बहती है; दाईं ओर मुड़कर वह उत्तर-पूर्व की ओर बढ़ती है: ये \omterm{def:g8:winds-global-circulation:global-winds}{पछुआ
पवनें} हैं।
\end{solution}

\begin{solution}{pb:g8:winds-global-circulation:1}
\textbf{1.} $59.6 - 23.5 = 36.1$ डिग्री \omterm{def:g6:latitude-longitude:longitude}{देशांतर}।

\textbf{2.} $36.1 \times 108 \approx 3899$: लगभग \qty{3900}{km}।

\textbf{3.} हाँ: $3900 - 3896 = 4$~km, एक नन्हा अंतर।

\textbf{4.} हाँ: दोनों \ang{14} उत्तर के पास हैं, \omterm{def:g3:continents-and-oceans:equator}{भूमध्य रेखा} और लगभग
\ang{30} उत्तर के बीच, उत्तर-पूर्वी व्यापारिक पवन की पेटी में।

\textbf{5.} उत्तर-पूर्व से (तीर दक्षिण-पश्चिम की ओर इशारा करते
हैं)।

\textbf{6.} आगे, और थोड़ा बगल की ओर: पवन जहाज़ के पीछे से, थोड़ा उसकी दाईं ओर से,
आती है और उसे पश्चिम की ओर धकेलती है।

\textbf{7.} यह \omterm{def:g8:winds-global-circulation:circulation}{वैश्विक परिसंचरण} का हिस्सा है: हवा \ang{30} उत्तर के पास लगातार नीचे
उतरती है और \omterm{def:g3:continents-and-oceans:equator}{भूमध्य रेखा} के पास ऊपर उठती है, इसलिए ज़मीनी हवा \omterm{def:g3:continents-and-oceans:equator}{भूमध्य रेखा} की ओर बहती
रहती है, और पृथ्वी का घूमना उसे हमेशा दाईं ओर मोड़ता
है।

\textbf{8.} व्यापारिक पवन जहाज़ के ख़िलाफ़ बहेगी। \omterm{def:g8:winds-global-circulation:global-winds}{पछुआ पवनें} ढूँढने के लिए कप्तान को
उत्तर की ओर, लगभग \ang{30} उत्तर से आगे, चलना चाहिए।

\textbf{9.} $7 \times 24 = \qty{168}{km}$ प्रति दिन।

\textbf{10.} $3 \times 168 = \qty{504}{km}$।

\textbf{11.} $3896 \div 168 \approx 23.2$ दिन।

\textbf{12.} लगभग \textbf{23 दिन}: एक महीने से कम।
\end{solution}
